\section{Numpy}\label{sec:06-numpy}


% -----------------------------
\subsection{Array erzeugen}\label{sec:06-array-erzeugen}
\begin{codebox}
\begin{lstlisting}[style=CodeExpert]
np.array([1, 2, 3])             # Array aus Python-Liste erstellen
np.asarray(x)                   # kein Copy falls x schon Array
np.copy(x)                      # tiefe Kopie: unabh. vom Original
np.zeros((3, 3))                # 3x3-Matrix, alle Werte = 0.0
np.ones((2, 4))                 # 2x4-Matrix, alle Werte = 1.0
np.full((2, 2), fill_value=9)   # 2x2-Matrix, alle Werte = 9
vk = np.full((gr1, gr2), fill_value=None)   # gr1xgr2-Array, alle Werte = None (dtype=object)
np.eye(4)                       # Einheitsmatrix: 1 auf Hauptdiag.
np.arange(1, 10, 2)             # [1,3,5,7,9] von 1 bis <10, Schr. 2
np.linspace(0, 1, 100)          # 100 gleichm. Werte von 0 bis 1
np.logspace(0, 3, 4)            # [1,10,100,1000]: 10^0 bis 10^3
np.empty((2, 2))                # nicht init. – zuf. Speicherinhalt
\end{lstlisting}
\end{codebox}


\subsection{Array-Informationen}\label{sec:06-array-informationen}
\begin{codebox}
\begin{lstlisting}[style=CodeExpert]
a.shape        # Tupel der Dimensionen (z.B. (3,4) = 3x4)
a.ndim         # Anzahl Dimensionen (2 fuer Matrix, 1 fuer Vektor)
a.size         # Gesamtanzahl Elemente (Produkt der Dimensionen)
a.dtype        # Datentyp der Elemente (z.B. int32, float64)
a.itemsize     # Bytes pro Element (z.B. 4 fuer int32)
a.nbytes       # Gesamtspeicher in Bytes (= size * itemsize)
a.T            # Transponierte Matrix
a.strides      # Schrittgroesse pro Achse in Bytes
a.flags        # Speicherlayout-Flags (C_CONTIGUOUS, etc.)
a.base         # Basisdaten: None=Owner, sonst Quelle des Views
\end{lstlisting}
\end{codebox}


\subsection{Typumwandlung}\label{sec:06-typumwandlung}
\begin{codebox}
\begin{lstlisting}[style=CodeExpert]
a.astype(np.float32)              # in float32 umwandeln (kopiert)
a.astype("int32", copy=False)     # kein Copy falls dtype schon passt
a.tolist()                        # zurueck zur Python-Liste
\end{lstlisting}
\end{codebox}

\subsection{Indexing, Slicing, Maskierung}\label{sec:06-indexing-slicing-maskierung}
\begin{codebox}
\begin{lstlisting}[style=CodeExpert]
a[1], a[1, :], a[:, 0]     # Element / ganze Zeile / ganze Spalte
a[1:4], a[::2]             # Slicing / jedes 2. Element
a[a > 3]                   # boolescher Indexzugriff (Maske)
np.where(a > 0, a, 0)      # bedingte Ersetzung: a wenn >0, sonst 0
np.nonzero(a)              # Indizes aller Nicht-Null-Elemente
np.take(a, [0, 2])         # Elemente per Index sammeln
np.put(a, [1, 3], [9, 10]) # Elemente per Index setzen (in-place)
np.compress(a > 2, a)      # Elemente filtern per boolescher Maske
\end{lstlisting}
\end{codebox}


\subsection{Broadcasting und Shapes}\label{sec:06-broadcasting-und-shapes}
\begin{codebox}
\begin{lstlisting}[style=CodeExpert]
a + b                                  # Shape automatisch erweitern
np.broadcast_shapes(a.shape, b.shape) # resultierenden Shape bestimmen
np.broadcast_to(a, (3, 4))            # manuelles Broadcasting (View)
np.expand_dims(a, axis=0)             # neue Achse einfuegen
np.squeeze(a)                         # Achsen der Groesse 1 entfernen
np.reshape(a, (2, -1))                # Form aendern (-1 = auto)
np.tile(a, (2, 3))                    # Muster 2x vertikal, 3x horiz.
\end{lstlisting}
\end{codebox}


\subsection{Stacking und Splitting}\label{sec:06-stacking-und-splitting}
\begin{codebox}
\begin{lstlisting}[style=CodeExpert]
# Kombinieren von Arrays
np.concatenate((a, b), axis=0)   # entlang bestehender Achse
np.stack([a, b], axis=0)         # neue Achse einfuegen
np.vstack([a, b])                # vertikal stapeln (Zeilen)
np.hstack([a, b])                # horizontal stapeln (Spalten)
np.dstack([a, b])                # tiefen-stapeln (3. Achse)
# Aufteilen von Arrays
np.split(a, 3)                   # in 3 gleiche Teile aufteilen
np.array_split(a, 3)             # in 3 Teile (ungleiche OK)
np.vsplit(a, 2)                  # vertikal in 2 Haelften teilen
np.hsplit(a, 2)                  # horizontal in 2 Haelften teilen
\end{lstlisting}
\end{codebox}


\subsection{Mathematik}\label{sec:06-mathematik}
\begin{codebox}
\begin{lstlisting}[style=CodeExpert]
np.add(a, b), np.subtract(a, b)      # Addition, Subtraktion
np.multiply(a, b), np.divide(a, b)   # Multiplikation, Division
np.mod(a, b), np.remainder(a, b)     # Modulo / Rest
np.exp(a), np.log(a), np.log10(a)    # e^a, ln(a), log10(a)
np.sqrt(a), np.square(a)             # Wurzel, Quadrat (a^2)
np.clip(a, min, max)                 # Werte auf [min, max] begrenzen
np.sign(a), np.abs(a)                # Vorzeichen (-1/0/1), Betrag
np.ceil(a), np.floor(a), np.round(a) # Aufrunden, Abrunden, Runden
\end{lstlisting}
\end{codebox}


\subsection{Aggregate und Statistik}\label{sec:06-aggregate-und-statistik}
\begin{codebox}
\begin{lstlisting}[style=CodeExpert]
np.sum(a), np.prod(a)          # Summe, Produkt aller Elemente
np.mean(a), np.median(a)       # arithm. Mittelwert, Median
np.std(a), np.var(a)           # Standardabweichung, Varianz
np.min(a), np.max(a)           # kleinster, groesster Wert
np.argmin(a), np.argmax(a)     # Index von Min bzw. Max
np.ptp(a)                      # Peak-to-Peak: max(a) - min(a)
np.percentile(a, 75)           # 75. Perzentil
np.histogram(a, bins=10)       # Histogramm berechnen
\end{lstlisting}
\end{codebox}


